\chapter{Реализация базы данных}

Реализация базы данных выполнена в PostgreSQL. Для удобства проверки был подготовлен файл \texttt{docker-compose.yaml}, который запускает контейнер PostgreSQL и автоматически выполняет SQL-файлы инициализации.

\section{Файлы проекта}

В проект входят следующие основные файлы:

\begin{enumerate}
    \item \texttt{schema.sql} создает структуру базы данных.
    \item \texttt{data.sql} добавляет тестовые данные.
    \item \texttt{queries.sql} содержит десять запросов для проверки работы базы.
    \item \texttt{diagram.png} содержит ER-диаграмму.
    \item \texttt{docker-compose.yaml} запускает PostgreSQL и подключает файлы инициализации.
    \item \texttt{README.md} содержит описание проекта и инструкцию по запуску.
\end{enumerate}

\section{Ограничения целостности}

В каждой таблице создан первичный ключ. Первичные ключи нужны, чтобы каждая запись имела уникальный идентификатор. В таблице \texttt{number\_assignments} используются внешние ключи на абонента, номер, SIM-карту, тариф и сотрудника. Это не дает создать выдачу, которая ссылается на несуществующие данные.

Ограничение \texttt{NOT NULL} используется для обязательных полей. Например, у абонента обязательно должно быть ФИО, у номера обязательно должен быть сам номер, а у SIM-карты обязательно должны быть ICCID и IMSI.

Ограничение \texttt{UNIQUE} используется для значений, которые не должны повторяться. Уникальными сделаны паспорт абонента, email, логин сотрудника, номер телефона, ICCID и IMSI. Это соответствует реальной логике работы оператора.

Ограничение \texttt{CHECK} используется для проверки статусов и числовых значений. Например, стоимость тарифа, количество минут, SMS и гигабайт не могут быть отрицательными. Статус номера может иметь только заранее заданные значения: \texttt{available}, \texttt{reserved}, \texttt{assigned} или \texttt{blocked}.

\section{Индексы}

В базе созданы индексы для полей, которые часто используются в поиске и фильтрации. Например, индекс по статусу номера помогает быстрее находить свободные или заблокированные номера. Индекс по времени операции нужен для отчетов за период. Индекс по абоненту в таблице выдач помогает быстрее получать все номера конкретного абонента.

Также созданы частичные уникальные индексы для активных выдач. Они не позволяют одному активному номеру или одной активной SIM-карте быть выданными двум абонентам одновременно. Это важное бизнес-правило для выбранной предметной области.

\section{Запуск через Docker}

Для запуска базы данных используется команда:

\begin{verbatim}
docker compose up -d
\end{verbatim}

При первом запуске контейнер PostgreSQL выполняет файлы \texttt{schema.sql} и \texttt{data.sql}. После этого база уже содержит структуру и тестовые данные. Для подключения можно использовать параметры:

\begin{verbatim}
host: localhost
port: 5432
database: sim_cards_db
user: sim_user
password: sim_password
\end{verbatim}

Для полной очистки данных можно выполнить:

\begin{verbatim}
docker compose down -v
\end{verbatim}

Эта команда удаляет контейнер и том с данными. После следующего запуска база будет создана заново из SQL-файлов.

\section{Тестовые данные}

Для проверки работы базы был подготовлен файл \texttt{data.sql}. В нем добавлены абоненты, тарифы, сотрудники, телефонные номера, SIM-карты, выдачи, блокировки и история операций. Всего в базе находится 75 тестовых записей. Этого достаточно, чтобы запросы возвращали не пустые результаты и показывали работу связей между таблицами.
